\documentclass[book.tex]{subfiles}
\begin{document}

\chapter{Electrolytes}
\label{chap:electrolytes}
\noindent\rule{11cm}{0.4pt} \\
\textbf{Prerequisites:} \autoref{chap:greens_identities}, \autoref{chap:maxwells_equations}, \autoref{chap:phase_equilibrium}  \\
\textbf{Difficulty Level:} ** \\
\noindent\rule{11cm}{0.4pt}
\vspace{5mm}

Any molecule that can dissolve in polar liquids, dissociate into ions, and carry an electrical current is an \textit{electrolyte}. Molecules that dissociate completely are called strong electrolytes, and those that dissociate only partly are called weak electrolytes. Such ions can regulate the chemical and physical equilibrium of other charged objects. For example, DNA is a molecular chain of negative phosphate charges. Because its charges repel each other, DNA is relatively expanded in water. When $NaCl$ is added to an aqueous solution of DNA, the salt dissociates into positive ions ($Na^+$) and negative ions ($Cl^-$). The positive ions seek the DNA's negative charges, surrounding and shielding them. This shielding weakens the intra-DNA charge repulsion, causing the DNA to become more compact. Two biological cells each having a net negative surface charge will also repel each other.

In this chapter we explore how ions shield charged objects. We start by a review of electrostatics where we derive the expression for the electrostatic energy of a charged sphere and cylinder. This expression will provide the starting point for the analysis of more complex situations, such as \textit{electrostatic screening} due to the large concentration of ions close to a surface with opposite charge.

\section{Electrostatic potential}

Electrostatic interactions in an aqueous environment dramatically influence the thermodynamic behavior. Polyelectrolytes (charged polymers) are soluble in water, in contrast to most hydrocarbon-based polymers. Most biopolymers are polyelectrolytes, and in many instances, their biological function hinges on the role of electrostatic interactions (Fig.~\ref{fig1}).

\begin{marginfigure}
\includegraphics[width=\linewidth]{figures/part2a/statistical_mechanics/electrolytes/fig1.png}
\caption{Hierarchical assembly of chromatin and the nucleosome core particle. Our meter-long genome is packaged into a micronsized nucleus with the aid of cationic proteins called histones.}
\label{fig1}
\end{marginfigure}

The force on a point particle with charge q in an electric field $\vec{E}$ is given by $\vec{f} = q\vec{E}$. The potential energy can then be written as $q\Psi$, where $\Psi$ is defined to be the electrostatic potential, giving $\vec{E}=-\nabla \Psi$.

The electric field satisfies \textit{Gauss' law}, which is written as:

\begin{align}\label{Eq1}
\int_S \vec{E}\cdot \vec{n} dA = \frac{4\pi}{\varepsilon}\int_V \rho dV
\end{align}

\noindent where $\varepsilon$ is the dielectric constant of the medium (for example, $\varepsilon = 7.083\times 10^{-10} C^2N^{-1}m^{-2}$, for water at room temperature), and $\rho$ is the local charge density.

We apply the \textit{divergence theorem} to this to arrive at:
\begin{align}\label{Eq2}
\int_V \vec{\nabla}\cdot \vec{E}dV = \frac{4\pi}{\varepsilon}\int_V \rho dV
\end{align}

Therefore, the electrostatic potential $\Psi$ is governed by Poisson's align:
\begin{align}\label{Eq3}
\nabla^2\Psi = -\frac{4\pi\rho(\vec{r})}{\varepsilon}
\end{align}

\section{Electrostatic potential around a charged sphere}

Consider a sphere with a radius a and total charge $q$. Due to symmetry, the electrostatic potential around the sphere is radial, \textit{\i.e.} $\Psi = \Psi(r)$. Poisson's equation outside of the sphere is written as:
\begin{align}\label{Eq4}
\frac{1}{r^2}\frac{d}{dr}r^2\frac{d\Psi}{dr} = 0 \rightarrow \Psi = \frac{A_1}{r} + A_2
\end{align}

\noindent where $A_1$ and $A_2$ are constant of integration (to be determined).

Assume as $r\rightarrow \infty$, $\Psi\rightarrow 0$; thus, $A_2 = 0$.

The electric field in the radial direction is $\hat{e}_r\cdot \vec{E} = \frac{A1}{r_2}$, and, using Gauss' law, we can write:
\begin{align}\label{Eq5}
\frac{A_1}{a^2}4\pi a^2 = \frac{4\pi q}{\varepsilon} \rightarrow A_1 = \frac{q}{\varepsilon} \rightarrow \Psi = \frac{q}{\varepsilon}\frac{1}{r}
\end{align}

This potential is valid outside the radius $a$ ($r > a$). For a point particle ($a\rightarrow 0$), this is valid throughout space.

The electrostatic potential provides the necessary information to
find the energy. The energy of interaction between two ions with charge $e$ (charge of an electron) separated by a distance $r$ is:
\begin{align}\label{Eq6}
E_{coulomb} = \frac{e^2}{\varepsilon}\frac{1}{r}
\end{align}

\noindent where $\varepsilon$ is the dielectric constant of the medium (for example, $\varepsilon = 7.083 \times 10^{-10}C^2N^{-1}m^{-2}$ for water at room temperature) and the electron charge $e = 1.6022\times 10^{-19}C$.

Define the Bjerrum length:
\begin{align}\label{Eq7}
l_B = \frac{e^2}{k_BT\varepsilon}
\end{align}

\noindent which gives the distance that two electrons are separated when their energy is equal to $k_BT$ ($l_B = 7\AA$ for water at room temperature).

The energy of interaction for two monovalent ions in a dielectric
medium is written as:
\begin{align}\label{Eq8}
E_{coulomb}= k_BT\frac{l_B}{r}
\end{align}

\section*{\textit{Electrostatic potential around a charged cylinder}}

Consider a negatively charged cylinder of radius $a$ with a surface
charge such that the average distance between the surface charges is $b$. The electrostatic potential in an uncharged medium ($\rho = 0$) is given by:
\begin{align}\label{Eq9}
\nabla^2\Psi = \frac{1}{r}\frac{\partial}{\partial r}r\frac{\partial}{\partial r}\Psi  = 0 \rightarrow \Psi = A_1\log r + A_2
\end{align}

\noindent where we express the Laplacian in cylindrical coordinates and assume $\Psi = \Psi(r)$ only.

We find the boundary condition on the cylinder surface using Gauss' law Using the divergence theorem, we write:
\begin{align}\label{Eq10}
\int_S \vec{E}\cdot \vec{n}dA = \frac{4\pi}{\varepsilon}\int_V \rho dV
\end{align}

This gives the boundary condition on the cylinder surface:
\begin{align}\label{Eq11}
\left( \hat{e}_r \cdot\vec{E} \right) 2\pi ab=-\frac{4\pi e}{\varepsilon}
\end{align}

This boundary condition on the cylinder surface is rewritten as:
\begin{align}\label{Eq12}
\frac{\partial \Psi}{\partial r}\bigg|_{r=a} = \frac{2e}{ab\varepsilon}
\end{align}

\noindent giving $A_1=\frac{2e}{b\varepsilon}$.

For simplicity, we set $\Psi\rightarrow 0$ as $r\rightarrow a$, fixing $A_2 = -A_1 \log a$. This gives the solution for the electrostatic potential.
\begin{align}\label{Eq13}
\Psi(r) = \frac{2k_BTq}{e}\log\left(r/a\right)
\end{align}

\noindent where $q = \l_B/b$ and $l_B = e^2/(k_BT\varepsilon)$.

Unlike the point charge, the electrostatic potential outside of a charged cylinder diverges logarithmically, which has dramatic consequences on the counter-ion distribution.

\section{Screened electrostatic interaction in salt (Debye-H{\"u}ckel theory)}

When salt is added to the solution, the nature of the electrostatic
interactions is dramatically altered. A charged molecule or surface attracts the mobile ions that have opposite charge, or \textit{counter-ions}, to the surface due to the strength of electrostatic interactions. At the same time, the charged surface repels the mobile ions of the same sign, called \textit{co-ions}. For example a polyanion (negative polyelectrolyte) has a flurry of positively charged counter-ions swarming around the polymer (Fig~\ref{fig2}). The interactions between charges are dramatically reduced due to the \textit{screening} of the electostatic interactions (Fig~\ref{fig3}).

\begin{marginfigure}
\includegraphics[width=\linewidth]{figures/part2a/statistical_mechanics/electrolytes/fig2.png}
\caption{Schematic of the salt-screening effect for a polyelectrolyte. The density of positive and negative counter-ions are colored red and blue, respectively. Electro-neutrality requires the total charge in solution to be zero, thus the bulk concentration is purple (equal positive and negative).}
\label{fig2}
\end{marginfigure}

Consider the interaction between two isolated ions in a monovalentsalt solution with concentration $n_s$ (for example, $NaCl$ that dissociates into $Na^+$ and $Cl^-$ ions). The electrostatic potential $\Psi$ around the tagged ion is governed by the Poisson's equation.
\begin{align}\label{Eq14}
\nabla^2\Psi = -\frac{4\pi\rho(\vec{r})}{\varepsilon}
\end{align}

\noindent where $\rho$ is the local charge density in the solution.

\begin{marginfigure}
\includegraphics[width=\linewidth]{figures/part2a/statistical_mechanics/electrolytes/fig3.png}
\caption{Schematic of the salt-screening effect between two isolated
ions.}
\label{fig3}
\end{marginfigure}

The energy of an ion with charge $q = ze$ (\textit{e.g.} $q=-e$ for $Cl^-$) located at $\vec{r}$ in the electrostatic potential $\Psi$ is given by $E_\Psi = ze\Psi(\vec{r})$.
As a charged colloidal particle is immersed on the electrolyte solution, a  gradient concentration of co-ions and conter-ions is developed from the surface. Assume the local concentration of counter-ions is given by a simple
Boltzmann weight of the ions in the electrostatic potential $\Psi$. Thus,
\begin{align}\label{Eq15}
\rho (\vec{r}) = n_s e \left[ \exp\left(-\frac{e\Psi(\vec{r})}{k_BT}\right) -\exp\left(+\frac{e\Psi(\vec{r})}{k_BT}\right)\right]
\end{align}

\noindent where the left term is due to local concentration of $Na^+$, and the right is due to $Cl^-$, and the position coordinate $\vec{r} = 0$ at the surface of the charged colloidal particle (Fig.~\ref{Poisson-Boltzmann})

\begin{marginfigure} \label{Poisson-Boltzmann}
\includegraphics[width=\linewidth]{figures/part2a/statistical_mechanics/electrolytes/Poisson-Boltzmann.png}
\caption{Schematic of the $Cl^-$ and $Na^+$ concentrations at a distance $r$ from a negatively charged surface at $r = 0$}
\label{fig3}
\end{marginfigure}

Therefore, the electrostatic potential is governed by:
\begin{align}\label{Eq16}
\nabla^2\Psi = \frac{8\pi n_s e}{\varepsilon}\sinh\left(\frac{e\Psi}{k_BT} \right)
\end{align}

This is the Poisson-Boltzmann equation. This nonlinear equation
does not have a simple, closed-form solution. Implicit in this theory is the neglect of local concentration fluctuations, akin to a mean-field approximation.

For sufficiently large lengths where $\Psi \rightarrow 0$, we can expand to lowest order.
\begin{align}\label{Eq17}
\nabla^2\Psi = \frac{8\pi n_s e^2}{\varepsilon k_B T}\Psi
\end{align}

We define the Debye length (Fig~\ref{fig4}):
\begin{align}\label{Eq18}
l_D = \sqrt{\frac{\varepsilon k_BT}{8\pi e^2n_s}}
\end{align}

\noindent and assume that $\Psi = \Psi(r)$ only.

We now have:
\begin{align}\label{Eq19}
\frac{1}{r^2}\frac{\partial }{\partial r^2}r^2\frac{\partial}{\partial r}\Psi = \frac{1}{l_D^2}\Psi
\end{align}

In the limit of large length separation $r$ and dilute solution of
counter-ions, we can write the electrostatic energy between two
ions in a salt solution as:
\begin{align}\label{Eq20}
E_{DH} = \frac{e^2}{4\pi \varepsilon}\frac{e^{-r/l_D}}{r} = k_BTl_B \frac{e^{-r/l_D}}{r}
\end{align}

This is the Debye-H{\"u}ckel theory of charge screening. The Debye-H{\"u}ckel theory adds a simple correction to the bare
electrostatic potential to account for counter-ion correlations.

\begin{marginfigure}
\includegraphics[width=\linewidth]{figures/part2a/statistical_mechanics/electrolytes/fig4.png}
\caption{Plot of the Debye length $l_D$ versus the salt concentration $n_s$ for monovalent salt in water at room temperature.}
\label{fig4}
\end{marginfigure}

In the absence of salt, the electrostatic interaction energy is extremely long range, scaling as $E_{coulomb} \sim 1/r$. Salt screening leads to a short-range electrostatic interaction energy
that decays exponentially with distance of separation.

\section{Charge behavior at short range (Manning condensation)}

Neglecting the nonlinear terms in the Poisson-Boltzmann equation
leads to a qualitatively incorrect behavior for distances very close to a charged object, \textit{e.g.} the surface of a polyelectrolyte.  For a dilute solution, the average distance between counter-ions
is $r_0 \sim n^{-1/3}_s$. At $r < r_0$, the mean-field approximation breaks down, and we must address the interactions explicitly.

Consider the polyelectrolyte chain to be a charged cylinder with
radius a (Fig.~\ref{fig5}), which we assume to be very small ($a\rightarrow 0$ in our analysis). The charged cylinder has negative charges smeared over its surface such that the average distance between negative charges is $\eta$. We want to find the behavior of an isolated ion interacting with the cylinder at distances less than the average inter-ion spacing $r_0$. 

\begin{marginfigure}
\includegraphics[width=\linewidth]{figures/part2a/statistical_mechanics/electrolytes/fig5.png}
\caption{Schematic for the short-range interaction between a charged cylinder (polymer) and a surrounding solution of dilute ions.}
\label{fig5}
\end{marginfigure}

Earlier in this chapter, we found the electrostatic potential around
a negatively charged cylinder in a dielectric medium to be:

\begin{align}\label{Eq21}
\Psi(r) = \frac{2k_BTq}{e}\log\left(r/a\right)
\end{align}

\noindent where $q = l_B/\eta$. The electrostatic energy between the cylinder and a positively charged monovalent ion is given by $E_{cyl}(r) = e\Psi(r)$.

The partition function for an isolated ion interacting with the charged cylinder within the distance r0 for small cylinder radius ($a \rightarrow 0$) scales as:
\begin{align}\label{Eq22}
Q&\sim \int_0^{r_0}  r \exp\left( -\frac{E_{cyl}(r)}{k_BT} \right) dr\\
&\sim \int_0^{r_0}r\exp(-2q\log(r)) dr\\
&\sim\int_0^{r_0}r^{1-2q}dr \\
&\sim \frac{1}{2-2q}r^{2-2q}\bigg|_0^{r_0}
\end{align}

The partition function diverges if q > 1.

To reconcile the divergence of the free energy, the mobile ions
condense on the cylinder surface to neutralize the surface charges such that effectively $q_{eff} = l_B/\eta_{eff} = 1$. At short length scales, Manning condensation leads to a reduced surface charge. At long length scales, the interaction is a screened electrostatic coulomb potential.

\end{document}
